\documentclass[11pt]{article} 
\usepackage{nopageno}
\setcounter{secnumdepth}{1} % Use 1 for sections, 2 for subsections, etc.
\usepackage[margin=1in]{geometry}
\begin{document}

% Remove page numbers---Research.gov will paginate the document instead
%
\pagestyle{empty} 

\noindent \textbf{Mentoring Plan:}

\noindent The project will involve: a PhD student advised by PI Napolitano and mentored by CoPI Chi through their committee; a post-doctoral scholar advised by CoPI Chi and mentored through regular meetings with PI Napolitano.
The mentoring plan will be updated annually based on the feedback and needs of the mentees. 
The success of the mentoring plan will be assessed through regular progress reports, participant feedback, and the participants’s academic and professional achievements. 
By providing mentoring and support, the project aims to develop well-rounded, competent, and confident researchers prepared to address complex challenges at the intersection of disaster science, historic preservation, social science, and data science. 
Mentoring activities will include but not be limited to: 

\noindent \textbf{Regular one-on-one meetings:} PI Napolitano will hold weekly meetings with the graduate student to discuss research progress, challenges, and opportunities. CoPI Chi will do the same with the post-doctoral scholar.

\noindent \textbf{Career counseling and professional development:} The mentees will be encouraged to present at conferences, workshops, and seminars. Additionally, PI Napolitano and CoPI Chi will provide career counseling to help the mentees identify their long-term career goals and develop strategies to achieve them. 
This will include discussions on career options within academia, industry, and government, as well as guidance on networking and job search strategies. 
The mentees will also have opportunities to network with industry partners and participate in career development events organized by PSU Career Services.

\noindent \textbf{Training in grant proposal and scholarly communication :} The mentee will receive training and guidance on preparing NSF grant proposals, peer-reviewed publications, and engaging presentations through the Graduate Writing Center. 
PI Napolitano and CoPI Chi will review drafts, provide constructive feedback, and help the mentees refine their writing and presentation skills. 
Mentees will also have access to the materials for the graduate-level research methods course that PI Napolitano teaches. 

\noindent \textbf{Teaching and mentoring skill development:} The mentees will have opportunities to assist in developing and delivering course content related to the project. 
PI Napolitano will provide guidance on effective teaching techniques, student engagement, and assessment and, as applicable, encourage them to pursue the Penn State Center for the Integration of Research, Teaching and Learning Certificate. 
The mentees will also receive mentorship on mentoring undergraduate students working on the project from the PSU Undergraduate Research and Fellowships Mentoring program.

\noindent \textbf{Collaboration and diversity training: }The mentees will be encouraged to collaborate with researchers from diverse backgrounds and disciplinary areas. 
PI Napolitano will provide guidance on effective communication, teamwork, and conflict-resolution skills. 
All students will also participate in diversity, equity, and inclusion trainings such as those offered at PSU on ``Addressing Unconscious Bias as a Leader" and ``Workplace Diversity, Equity, Inclusion, and Belonging in Action." 

\noindent \textbf{Responsible professional practices:} Both mentees will receive training in responsible conduct of research, including research ethics, data management, and intellectual property rights through the PSU Scholarship and Research Integrity Education program. PI Napolitano and CoPI Chi will also discuss professional practices, such as time management, work-life balance, and professional etiquette.

\end{document}